\reportsection{The Parametrisation Model}
\label{sec:spine}

Binding and hiding failed separately and were sized separately, but the two
sizings are one rule, and that they share a rule is the point of this section.
The earlier \textsf{Arkworks/crypto-primitives}~\cite{ArkCryptoPrimitives} bought security by hardcoding every
choice, and developers, wanting simpler hash interfaces, variable tree shapes,
and efficient multi-opening compression, asked for that freedom back.  Granting
it lets a user control $(\field,\secpar,\RO,\text{interface})$ in full, and the
danger is exactly that without hardcoded guardrails a misconfiguration stays
silent.  So the developer faces a question they cannot answer from the interface
alone: which security claim does my concrete configuration actually support?
Answering it is our obligation, not theirs.  Given those choices and a
declared adversary model, the
infrastructure owes back the parameters that must be used, a proof that they
meet the target, the cost they incur, and the property they buy.  The user can
then keep their attention on their own implementation and read its concrete
security off the guarantees we give.

\subsection{The Object It Sizes}
\label{sec:object}

We fix the commitment object once, following the Chiesa--Yogev notation so the
mathematical boundary does not drift from the reference the implementation
follows~\cite{CY26}.

\begin{definition}[Vector commitment scheme]\label{def:vcgeneral}
A \emph{vector commitment scheme}~\cite{CF13} over alphabet $\MTAlphabet$ and length
$\MTMessageLength$ is a triple $\VC=(\Commit,\Open,\VCCheck)$:
\begin{itemize}
  \item $\Commit(\msgvec)\to(\comm,\aux)$ binds a message
  $\msgvec\in\MTAlphabet^{\MTMessageLength}$ to a short commitment $\comm$ and
  retains committer state $\aux$;
  \item $\Open(\aux,\idxset)\to\opening$ authenticates an index set
  $\idxset\subseteq[\MTMessageLength]$;
  \item $\VCCheck(\comm,\idxset,\MTMessageSubVector,\opening)\to\bits$ decides
  whether $\opening$ authenticates the subvector $\MTMessageSubVector$ at
  $\idxset$ against $\comm$, where for an honest committer
  $\MTMessageSubVector$ is exactly the committed message $\msgvec$ restricted
  to the positions in $\idxset$.
\end{itemize}
A scheme is \emph{complete} (\cref{def:vc-complete}) and
\emph{position-binding} (\cref{def:vc-binding}), and may implement optional
extensions---hiding, treated here, plus updatability, aggregation, functional
opening, and knowledge soundness---which \textsf{ark-vc}'s capability model
accommodates but this report does not treat.
\end{definition}

\begin{definition}[Completeness]\label{def:vc-complete}
$\VC$ is \emph{complete} if for every $\msgvec\in\MTAlphabet^{\MTMessageLength}$
and every $\idxset\subseteq[\MTMessageLength]$ the honest
$\Commit$--$\Open$--$\VCCheck$ round trip accepts $\msgvec_{\idxset}$ with
probability~$1$.  For Merkle this is honest reconstruction of the root from an
opening.
\end{definition}

\begin{definition}[Merkle commitment scheme]\label{def:vc}
The \emph{Merkle commitment scheme}~\cite{Merkle88}
$\MTSymbol=(\MTCommit,\allowbreak\MTOpen,\allowbreak\MTCheck)$
instantiates \cref{def:vcgeneral} in the random-oracle model over a general
Merkle shape $G_\ell$ of depth $d=\max_i d_i$, the deepest of the per-leaf depths
$d_i$ over leaves $i$: the three algorithms query a
shared oracle $\ROFunction$, the commitment is the tree root, and the committer
state is the retained salts and labels.  The binary tree is the arity-$2$ subset
of this shape.  Two parameters carry the security: the digest
width $\digestbits$ governs binding and the salt size $\MTSaltSize$ governs
hiding, with no tuning of one supplying what the other lacks.
\end{definition}

\begin{limitation}[No trapdoor, so no standard-model equivocation]
\label{lim:no-trapdoor}
$\MTSymbol$ has no cryptographic $\Setup$: the public parameter is the hash
itself, so the trapdoor is empty ($\td=\bot$).  This shapes what can even be
asked of the scheme.  Extraction cannot program a trapdoor and must instead
argue from the adversary's random-oracle query trace, and standard-model
equivocation, which opens one commitment to two values precisely by
programming a trapdoor, has nothing to program.  Equivocation is therefore
unreachable for this scheme by construction, not merely unproved.
\end{limitation}

\subsection{The Rule}
\label{sec:rule}

The rule itself was stated in the box of \cref{sec:background}: from the
designer's interface-level inputs $(\field,\secpar,H,\text{interface})$ and
the declared adversary budget $h$, the constructor
derives the digest width $D$, the field-salt cardinality $k$, and the byte-salt
cardinality $S$ of \cref{eq:rule}, and either meets each obligation or refuses
the configuration.  What remains to establish here is what the rule does not
depend on, and what it produces per field.

The rule is arity-independent.  $(D,k,S)$ depends only on $(\field,\secpar,h)$,
while the tree shape and arity enter only through the query budget $h$ and the
cost.  The same target
therefore sizes any Merkle shape and produces different outputs per field: for
BabyBear ($b=30$, $h=64$) it gives $D=9$, $k=5$, $S=16$; for BN254
($b\approx254$) it gives $D=2$, $k=1$, $S=16$.  These are not interchangeable,
and the $44$-second break of \cref{sec:background} used a $D=1$ configuration,
one of the hashers \textsf{Arkworks} currently ships, which is exactly the kind
of undersizing the constructor now refuses.
